\section{Comparison with the proposal archive and verification}\label{sec:verify}
The comparison uses ProveIt commit \href{https://github.com/VladimirReshetnikov/ProveIt/commit/89d12dd40424f1217fd818d84877b3da52396a3b}{\texttt{89d12dd4}}. Its \src{Main}{5}{readable Lean project} has Git tree \texttt{2e133df6}, identical to the readable project in upstream commit \href{https://github.com/lsernau/grh-lean/tree/72fead102f663d5c7a407964a13dc0cf9d977207/grh}{\texttt{72fead10}}. The \docref{guide}{human-readable guide} traces the digit construction in Section 8, the finite-field labels in Section 10, and the local operator assembly in Section 14. These are the three passages for which the present note supplies substitutions.

\begin{center}\small
\begin{tabularx}{\linewidth}{@{}p{.30\linewidth}X@{}}\toprule
Prior proposal, linked to source & Substitutions already proposed\\\midrule
\docref{proposals/centered_theta_lift_simplification}{Centered theta lifts} & Centered rank increment, ramified-source spanning, adaptive completion of a nearly full basis, and a random-product variant.\\[3pt]
\docref{proposals/proof_simplifications}{Three shorter steps} & Two-polynomial and Fourier estimates for the auxiliary prime, a Perron argument at the far cutoff, and a stronger spanning estimate.\\[3pt]
\docref{proposals/rh-simplifications}{Simplifications of RH} & Rank recurrence on the ramified source, simultaneous torsion interpolation for choosing twists, and a quadratic-large-sieve selector.\\[3pt]
\docref{proposals/rh_sparse_theta_proof}{Sparse cutoffs} & Reciprocal-convex interpolation, sparse or averaged cutoffs, and a revised quadratic-twist selection argument.\\[3pt]
\docref{proposals/shorter-theta-lift-proof}{Shortening the theta-lift proof} & Smooth blocks and a large-sieve route for choosing the quadratic field, ramified-source spanning, and endpoint interpolation.\\\bottomrule
\end{tabularx}\end{center}

In particular, torsion interpolation chooses a character before constructing the one-generator operators. Section~\ref{sec:flag} combines operators after that arithmetic work is finished. Adaptive basis completion estimates the cost of missing target vectors in the covolume comparison. The flag used here is inside one local lift space and its successive quotients all have dimension one. In Section~\ref{sec:denser-labels} the count concerns local polynomial types before the spanning argument starts. In Section~\ref{sec:paraboloid} the sparse objects are infinite places of a CM field at one fixed cutoff. The sparse-cutoff proposal varies the cutoff itself. These distinctions allow the three changes to be evaluated even after the existing proposals have been implemented.

\paragraph{Verification of the baseline.}
The \href{https://math.mit.edu/~drew/shortgrh/}{independent verification record of Andrew Sutherland} reports a full build, kernel replay of every module, and a successful paranoid Comparator run for the original readable edition. Its toolchain is Lean 4.35.0 with Mathlib revision \texttt{56a00585}. The equality of Git trees identifies which source those checks cover. This report does not reproduce that build.

\paragraph{Verification of the replacements.}
The attached standalone Lean files use Lean 4.32.0 and its standard library. \path{formal/IntegralPivots.lean} proves the existence of each successive pivot, the extension of the flag by a normalized integral vector, and the equivalence between coordinate integrality and integrality of its expansion. \path{formal/FlagRows.lean} proves in all finite dimensions that selecting the prescribed rows preserves clearing and gives the stated determinant. It defines the determinant by cofactor expansion and proves the triangular formula for that definition. \path{formal/DigitRigidity.lean} proves the modular rigidity of the explicit code for all admissible finite parameters, including the range bound, decoding, and injectivity. \path{formal/LabelArithmetic.lean} proves the integer obstruction to orbit collapse, invariance of the residue tests under the corresponding shifts, and the prime-power divisibility used in the norm estimate. The new modules are separate from the pinned project. The paper supplies the identification of the integral evaluation lattice with adelic forms, the finite-field orbit count and its application to the projector, and the transfer of the digit code to Frobenius orbits of places. Integration and measurement of the revised project remain open.

The three exact-arithmetic programs exercise the pivot construction and clearing on rational models at several prime valuations, the orbit and norm tests in finite fields, and the digit and signed-place conditions. The archive includes their outputs, the Lean axiom reports, and a source index with immutable links. Its checker verifies the indicated lines and Git blob hashes against a checkout. The finite examples supplement the uniform proofs.

\paragraph{Which change to implement first.}
The row construction can be inserted between the existing one-generator and determinant lemmas, then reused for the compressed determinant. The digit code can be integrated separately by strengthening only the eventual hypotheses on its fixed dimension. The projector change has the largest quantitative effect among the three local estimates and the widest integration surface. The test for success there is to build both the trace-comparison and Galois branches for the enlarged projector, as well as the final bad-type estimate. No global improvement to the required field degree or to the running time follows solely from the new local threshold.

\paragraph{Further questions.}
The individual bounds $B_j$ can be summed before taking a worst case. Do the Selmer estimates yield a smaller sum for the families actually used in the proof? The finite-field criterion also leaves room for an optimum. Roots whose exponent is divisible by $3$ can still retain their full orbit after every prime-power map if the remaining order in the quotient by $\mathbb F_{q^{r/3}}^\times$ has two distinct prime divisors. Counting these additional labels would require information about $H$, which the uniform estimate here avoids.
